% methodology.tex -- Enterprise Node Orchestrator Paper (RU-V5)

\section{Methodology}

\subsection{State Machine Design}

The orchestrator is modeled as a typed finite state machine
$\mathcal{M} = (S, E, \delta, s_0)$ where $S$ is the set of states,
$E$ is the set of events, $\delta: S \times E \to S$ is the transition
function, and $s_0 = \text{Idle}$ is the initial state.

\subsubsection{States}

The state set $S = \{\text{Idle}, \text{Receiving}, \text{Batching},
\text{Proving}, \text{Submitting}, \text{Error}\}$ captures the six
phases of the validium proof cycle:

\begin{itemize}
\item \textbf{Idle}: No pending work. The node polls for new transactions.
\item \textbf{Receiving}: Actively accepting transactions from enterprise
  systems via REST or WebSocket API.
\item \textbf{Batching}: Batch threshold reached. The node forms a batch,
  applies state updates to the SMT, and generates the circuit witness.
\item \textbf{Proving}: ZK proof generation in progress. Runs in a child
  process to avoid blocking the event loop.
\item \textbf{Submitting}: Proof generated. DAC attestation and L1
  settlement in progress.
\item \textbf{Error}: Recoverable error state. Transitions back to Idle
  via retry with exponential backoff.
\end{itemize}

\subsubsection{Transition Table}

Table~\ref{tab:transitions} enumerates all 17 valid transitions. The
critical design feature is pipelined receiving: the \textsc{TransactionReceived}
event is accepted in the Proving and Submitting states, enabling continuous
ingestion while a proof cycle completes.

\begin{table}[htbp]
\centering
\caption{State Machine Transition Table}
\label{tab:transitions}
\begin{tabular}{lll}
\toprule
\textbf{Source} & \textbf{Event} & \textbf{Target} \\
\midrule
Idle       & TransactionReceived   & Receiving \\
Idle       & ShutdownRequested     & Idle \\
\midrule
Receiving  & TransactionReceived   & Receiving \\
Receiving  & BatchThresholdReached & Batching \\
Receiving  & ErrorOccurred         & Error \\
\midrule
Batching   & WitnessGenerated      & Proving \\
Batching   & ErrorOccurred         & Error \\
\midrule
Proving    & ProofGenerated        & Submitting \\
Proving    & TransactionReceived   & Proving \\
Proving    & ErrorOccurred         & Error \\
\midrule
Submitting & BatchSubmitted        & Submitting \\
Submitting & L1Confirmed           & Idle \\
Submitting & TransactionReceived   & Submitting \\
Submitting & ErrorOccurred         & Error \\
\midrule
Error      & RetryRequested        & Idle \\
Error      & ShutdownRequested     & Idle \\
\bottomrule
\end{tabular}
\end{table}

\subsubsection{Invariants}

Six invariants govern correctness and are formalized for downstream
TLA+ specification:

\noindent\textbf{INV-NO1} \emph{Liveness}: If $|\text{pendingTx}| > 0$ and
$s = \text{Idle}$, then eventually $s = \text{Proving}$.

\smallskip\noindent\textbf{INV-NO2} \emph{Safety}: If $s = \text{Submitting}$, then
$\text{proof.prevRoot} = \text{smt.prevRoot}$ and
$\text{proof.newRoot} = \text{smt.currentRoot}$.

\smallskip\noindent\textbf{INV-NO3} \emph{Privacy}: The only data transmitted outside
the node boundary are: proof $(a, b, c)$, public signals
$(r_\text{prev}, r_\text{new}, \text{batchNum}, \text{enterpriseId})$,
and Shamir shares to DAC nodes.

\smallskip\noindent\textbf{INV-NO4} \emph{Crash Recovery}: After crash and restart,
$|\text{walReplayed}| + |\text{committed}| = |\text{totalEnqueued}|$.

\smallskip\noindent\textbf{INV-NO5} \emph{State Root Continuity}: For batch $N$,
$\text{submit}(N).\text{prevRoot} = \text{batch}(N-1).\text{newRoot}$.

\smallskip\noindent\textbf{INV-NO6} \emph{Single Writer}: Only the batch loop
modifies the SMT. No concurrent write access is permitted.

\subsection{Pipelined Architecture}

The sequential execution model (Idle $\to$ Receive $\to$ Batch $\to$
Prove $\to$ Submit $\to$ Idle) blocks transaction ingestion during proof
generation, creating unavailability windows of 12+ seconds. Following
the zkSync Era~\cite{zksync2025lifecycle} and push0~\cite{push0_2025orchestration}
patterns, we decompose the pipeline into three concurrent loops:

\begin{enumerate}
\item \textbf{Ingestion Loop}: Continuously accepts transactions via the
  Fastify REST/WebSocket API, writes to the WAL, and enqueues to the
  transaction queue.

\item \textbf{Batch Loop}: Monitors the queue size. When the batch
  threshold is reached, dequeues transactions, applies SMT updates,
  and generates the circuit witness.

\item \textbf{Proving/Submission Loop}: Takes formed batches, invokes
  the ZK prover (child process), distributes witness data to the DAC,
  and submits the proof to L1.
\end{enumerate}

The pipelined speedup arises when the batch loop prepares the next
batch while the proving loop processes the current one. The speedup
is bounded by:

\begin{equation}
\text{Speedup} = \frac{N \cdot T_\text{seq}}{T_\text{cold} + (N-1) \cdot \max(T_\text{prep}, T_\text{prove+submit})}
\end{equation}

\noindent where $T_\text{seq}$ is the sequential batch time, $T_\text{cold}$
is the first batch (no overlap), $T_\text{prep}$ is preparation time
(batch formation + witness generation), and $T_\text{prove+submit}$ is
the proving and submission phase. The speedup increases with batch size
because $T_\text{prep}$ grows (more witness generation work) while
$T_\text{prove+submit}$ grows more slowly.

\subsection{Component Integration}

The orchestrator integrates five components, each characterized by prior
research units in this pipeline:

\begin{table}[htbp]
\centering
\caption{Component Latencies from Prior Research}
\label{tab:components}
\begin{tabular}{llrl}
\toprule
\textbf{Component} & \textbf{Operation} & \textbf{Latency} & \textbf{Source} \\
\midrule
SparseMerkleTree & Insert (d32) & 1.825~ms & RU-V1 \\
SparseMerkleTree & Proof generation & 0.018~ms & RU-V1 \\
TransactionQueue & WAL write (group) & 149~$\mu$s & RU-V4 \\
BatchAggregator  & Batch formation & 0.02~ms & RU-V4 \\
BatchBuilder     & Witness gen (d32, b8) & 578~ms & RU-V2 \\
DACProtocol      & Full attestation & 163~ms & RU-V6 \\
StateCommitment  & L1 submission & 268K gas & RU-V3 \\
\bottomrule
\end{tabular}
\end{table}

\subsection{API Design}

The node exposes a Fastify-based REST API and WebSocket event stream for
enterprise client integration:

\begin{itemize}
\item \texttt{POST /v1/transactions} -- Submit a single transaction
\item \texttt{POST /v1/transactions/batch} -- Submit multiple transactions
\item \texttt{GET /v1/status} -- Node health and current state
\item \texttt{GET /v1/batches/:id} -- Query batch status and proof
\item \texttt{GET /v1/state/:enterprise} -- Current state root
\end{itemize}

WebSocket events (\texttt{ws://node/v1/events}) provide real-time
notifications for transaction acceptance, batch proving progress, and
L1 confirmation.

\subsection{Technology Selection}

Fastify v4 was selected over Express.js based on published
benchmarks~\cite{fastify2025comparison} showing 2--4$\times$ higher
request throughput with native TypeScript support and JSON schema
validation. WebSocket integration uses the
\texttt{@fastify/websocket}~\cite{fastify2025websocket} plugin for
real-time enterprise event streaming.
